\section*{Homogénéisation de l'équation de Poisson}
À titre d'extension, on présente également l'homogénéisation de 
l'équation de Poisson, qui permet de décrire explicitement les 
effets de charge d'espace lorsque l'hypothèse d'électroneutralité n'est plus vérifiée.

On réutilise ici, pour les fluctuations du potentiel $\widetilde{\varphi}_l$,
la fermeture introduite pour le terme migratoire. Cette analogie constitue
toutefois une approximation plus forte encore dans le cas de Poisson :
le problème de fermeture y est en principe couplé aux fluctuations de
concentration, puisque c'est précisément la distribution de charge qui
source le potentiel. Traiter ce problème de fermeture de façon découplée,
comme pour la diffusion ou la migration, revient à négliger ce couplage. La
dérivation qui suit doit donc être comprise comme indicative, fournissant la
structure du modèle homogénéisé plutôt qu'une fermeture rigoureuse.

En appliquant la moyenne superficielle à l'équation de Poisson :

\begin{equation}
-\varepsilon_0 \varepsilon_r \avg{\nab^2 \varphi_l} = F\sum_i z_i \avg{C_i}
\end{equation}

\subsubsection*{Membre de gauche}

En appliquant le théorème de la moyenne spatiale \eqref{eq:moy_spatiale} :

\begin{equation}
\avg{\nab\nab\varphi_l} = \nab\avg{\nab\varphi_l} + \frac{1}{V}\int_{A_{sf}} \mathbf{n}\,\nab\varphi_l\, dA
\end{equation}

La condition limite à l'interface $A_{sf}$ traduit la présence d'une 
charge de surface $\sigma_q$ :

\begin{equation}
-\varepsilon_0 \varepsilon_r \nab\varphi_l \cdot \mathbf{n} = \sigma_q 
\quad \text{sur } A_{sf}
\quad \Rightarrow \quad 
\mathbf{n}\cdot\nab\varphi_l = -\frac{\sigma_q}{\varepsilon_0 \varepsilon_r}
\end{equation}

En utilisant la fermeture introduite précédemment pour
$\widetilde{\varphi}_l$, tout en gardant à l'esprit qu'elle constitue
une approximation en présence du couplage avec l'équation de Poisson,
on obtient :
\begin{align}
\avg{\nab\nab\varphi_l}
&= \nab\avg{\nab\varphi_l} - \frac{a_v \sigma_q}{\varepsilon_0 \varepsilon_r} \notag\\
&= \nab\!\left(\nab\avg{\varphi_l}^l + \frac{1}{V}\int_{A_{sf}} \mathbf{n}\,\varphi_l\, dA\right) 
- \frac{a_v \sigma_q}{\varepsilon_0 \varepsilon_r} \notag\\
&= -\frac{a_v \sigma_q}{\varepsilon_0 \varepsilon_r}
+ \nab\!\left[\varepsilon_l \nab\avg{\varphi_l}^l
+ \frac{1}{V}\int_{A_{sf}} \mathbf{n}\,\widetilde{\varphi}_l\, dA\right] \notag\\
&= -\frac{a_v \sigma_q}{\varepsilon_0 \varepsilon_r}
+ \nab\!\left[\varepsilon_l \nab\avg{\varphi_l}^l
\left(I + \frac{1}{V_l}\int_{A_{sf}} \mathbf{n}\,\mathbf{d}\, dA\right)\right] \notag\\
&= -\frac{a_v \sigma_q}{\varepsilon_0 \varepsilon_r}
+ \nab\!\left[\varepsilon_l \frac{\varepsilon^{\text{eff}}}{\varepsilon_r} \nab\avg{\varphi_l}^l\right]
\end{align}

où $\varepsilon^{\text{eff}}$ est la permittivité effective du milieu poreux, 
définie de façon analogue à $D_{\text{eff}}$ :

\begin{equation}
\varepsilon^{\text{eff}} = \varepsilon_r \left(I + \frac{1}{V_l}\int_{A_{sf}} \mathbf{n}\,\mathbf{d}\, dA\right)
\end{equation}

\subsubsection*{Membre de droite}

En utilisant la relation entre moyenne superficielle et moyenne intrinsèque 
\eqref{eq:relation_moy} :

\begin{equation}
F\sum_i z_i \avg{C_i} = F\sum_i z_i \varepsilon_l \avg{C_i}^l
\end{equation}

\subsubsection*{Réunification}

En combinant les membres de gauche et de droite, on obtient l'équation 
de Poisson homogénéisée :

\begin{equation}
\boxed{
-\varepsilon_0 \nab\!\left(\varepsilon_l \varepsilon^{\text{eff}} \nab\avg{\varphi_l}^l\right)
= F\sum_i z_i \varepsilon_l \avg{C_i}^l + a_v \sigma_q
}
\label{eq:poisson_homo}
\end{equation}
